%-------------------------------------------------------------------------------
%  SECTION TITLE
%-------------------------------------------------------------------------------
\cvsection{\'Education}

%-------------------------------------------------------------------------------
%  CONTENT
%-------------------------------------------------------------------------------
\begin{cventries}

  %---------------------------------------------------------
  \cventry
  {Université de Montpellier} % Institution
  {Doctorat en Robotique Humanoïde et Vision par Ordinateur} % Degree
  {Montpellier, France} % Location
  {Oct. 2014 -- Nov. 2018} % Date(s)
  {
    \begin{cvitems}
    \item \entrypositionstyle{Titre} : « SLAM visuel pour la localisation et la commande en boucle fermée de
      robots humanoïdes » (Soutenue en novembre 2018).
    \item \entrypositionstyle{Cadre} : Thèse opérée en co-direction LIRMM / I3S / AIST-JRL.
    \end{cvitems}
  }
  %---------------------------------------------------------

  %---------------------------------------------------------
  \cventry
  {Polytech Nice Sophia -- Université de Nice} % Institution
  {Diplôme d'Ingénieur en Informatique (Grade de Master)} % Degree
  {Sophia Antipolis, France} % Location
  {Sept. 2011 -- Sept. 2014} % Date(s)
  {
    \begin{cvitems} % Description(s) bullet points
    \item \entrypositionstyle{Spécialisation} : Vision, Image et Multimédia.
    \item \entrypositionstyle{Année ERASMUS au Trinity College Dublin (2012-2013)} :
      \begin{itemize}
        \item \emph{Calcul \& Simulation (C++)} : conception de moteurs physiques temps réel (simulation de
          fluides, collisions) au sein du Master des technologies interactives.
        \item \emph{Stage au CRANN Institute, Dublin (3 mois)}: Conception d'un logiciel graphique (C++, Qt)
          de recalage interactif de courbes pour la recherche en microscopie à effet tunnel.
      \end{itemize}
    \item \entrypositionstyle{Réalité Virtuelle \& Augmentée} : Premier projet de recherche supervisé par
      \emph{Dr. Andrew Comport} sur le rendu immersif de cartes SLAM 3D dans un casque de réalité virtuelle.
    \item \entrypositionstyle{Stage au Technische Universität München (Allemagne, 2014)} : Projet de fin
      d'études (6 mois) sur la reconnaissance de lieux par réseaux de neurones convolutionnels (CNN)
      supervisé par \emph{Dr. Jurgen Sturm} dans la chaire de \emph{Prof. Daniel Cremers}
    \item \entrypositionstyle{Engagement Sociétal} : Co-développement d'un logiciel applicatif 3D adapté
      spécifiquement pour les personnes malvoyantes.
    \end{cvitems}
  }

  %---------------------------------------------------------
  \cventry
  {Lycée Kérichen} % Institution
  {Classes Préparatoires aux Grandes Écoles (CPGE MPSI)} % Degree
  {Brest, France} % Location
  {Sept. 2009 -- Juin 2011} % Date(s)
  {
    % \begin{cvitems} % Description(s) bullet points
    %   \item \entrypositionstyle{Filière} : Mathématiques, Physique et Sciences de l'Ingénieur (MPSI / MP).
    % \end{cvitems}
  }

  %---------------------------------------------------------
  \cventry
  {Lycée de l'Harteloire} % Institution
  {Baccalauréat Scientifique} % Degree
  {Brest, France} % Location (Ajuste la ville si besoin)
  {Juin 2009} % Date(s)
  {
    % \begin{cvitems} % Description(s) bullet points
    %   \item \entrypositionstyle{Filière} : Scientifique
    % \end{cvitems}
  }
\end{cventries}
